\section{\label{II_2_B}Entraînement}

L'analyse automatique du contenu des images repose donc sur les architectures des \gls{cnn}, fondées sur des méthodes d'\gls{Apprentissage profond}, dont YOLO est un exemple concret. Toutefois, les algorithmes fondés sur l'\gls{Apprentissage profond} subissent généralement un entraînement supplémentaire afin d'analyse adéquatement le contenu des images en fonction d'un contexte donné. Deux approches méthologiques d'entraînement peuvent être appliquées, l'\gls{appsup} et l'\gls{appnonsup}. L'\gls{appsup} est fondé sur l'utilisation de données annotées pour l'amélioration des prédictions du modèle d'analyse d'image, tandis que l'\gls{appnonsup} exécute des tâches d'analyse et de classification à partir de données brutes, non annotées \footcite{delua2024}. Toutefois, l'\gls{appnonsup} ne permet pas d'adapter suffisamment le modèle de détection d'objets au corpus d'images spécifique étudié. En effet, si Yolov8 a toutefois fait l'objet d'un pré-entraînement en amont de sa mise à disposition, sur le jeu de données COCO, ce dernier, composé de 200 000 images annotées, reste généraliste et adapté à la détection d'objets communs et contemporain\footcite{2023b}. Or, l'efficacité des modèles basés sur l'\gls{Apprentissage profond} sur un corpus spécifique dépend fondamentalement de nature et de la qualité de leur données d'entraînement. 
Au sein des sous sections suivantes, nous prendrons comme exemple d’application le projet des archives britanniques en ligne, qui propose une chaine de traitement pour la détection d’objet basée sur l'\gls{Apprentissage profond} et la supervision humaine \footcite{kasturi2024}. Le projet, à travers la mise en place d'une supervision au sein du processus d'analyse par le modèle, propose de mettre en application le concept de \gls{hil} \footcite{kasturi2024}.
Le corpus étudié, caractérisé par l'ancienneté de ses images en fait un corpus d’images spécifique sur lequel l’algorithme doit subir un entraînement supplémentaire, supervisé. Afin d'adapter le modèle à notre corpus et ainsi d'en améliorer l'efficacité, une phase de pré traitement est nécessaire sur les données du corpus, au cours de laquelle des images sont sélectionnées afin d’être annotées pour être pré-testées par le modèle \footcite{kasturi2024}. La supervision humaine est déterminante dans ce processus, particulièrement dans le cas d’analyse de corpus inédits sur lesquels l’algorithme reste à entraîner \footcite{kasturi2024}. 
Les british online archives implémentent l'adaptation du modèle par la pré selection d'un nombre d'images définis, annotées manuellement à l'aide de la plateforme \gls{labels} \footcite{kasturi2024}. \gls{labels} est une plateforme open source déployable localement, qui permet ainsi de respecter la protection des données analysées \footcite{zotero-item-783}.
Afin d'éviter l'annotation manuelle de l'ensemble des \gls{imgcle} issues du processus de segmentation, il peut être envisagé d'exécuter au préalable \gls{yolo} depuis \gls{labels} sur l'ensemble des \gls{imgcle} du corpus afin de générer une base de prédictions \footcite{development2025}. Les détections brutes générées peuvent ensuite être importées au sein de la plateforme comme une couche d'annotation préalable à partir de laquelle affiner le modèle. Au sein de \gls{labels}, les prédictions s'affichent sous la forme de \textit{bounding boxes}, c'est-à-dire de boîtes englobantes autour des objets detectés dans l'image, avec le score de confiance qui y est associé \footcite{kasturi2024}. Le paramètre [choices] permet au modèle de proposer une classification sur les objets détectés, c'est-à-dire un mot clé \footcite{zotero-item-785}. Afin de focaliser la supervision du modèle sur les éléments les plus pertinents, il peut être décidé de ré annoter les images dont la détection est jugée incertaine par le modèle via un score de confiance en dessous du seuil minimal attendu \footcite{development2025}. Les archives brittaniques préconisent également la uppression des boites englobantes redondantes afin de sélectionner celles qui sont les plus pertinentes pour la supervision \footcite{kasturi2024}. Dans \gls{labels}, les utilisateurs (une équipe technique dans le cas de Skinsoft) examinent les prédictions ainsi générées sur les images, rectifient si besoin la définition des "boites" et la classification associée \footcite{kasturi2024}. Les données ainsi corrigées sont exportées au format \gls{yolo} depuis \gls{labels}, puis réinjectées dans la chaîne de réentraînement du modèle afin d'améliorer les performances en accord avec le corpus \footcite{kasturi2024}.
Cette méthode s'appuie sur le transfert d’apprentissage (transfert learning), qui consiste à \textit{fine tuner} un modèle existant, c’est-à-dire de l’affiner, afin de l'adapter au corpus analysé \footcite{bhargav2019}. À défaut de réentraîner le modèle dans sa globalité, les prédictions réalisées à partir du jeu de données d'entraînement COCO sont conservées puis spécialisées sur un corpus donné \footcite{bhargav2019}. L'affinage du modèle via la rectification et l'annotation des données, constitue une étape de prétraitement, en amont de l’analyse des données, afin d'améliorer l'efficacité de la détection d’objets \footcite{kasturi2024}.
Cette étape de pré-traitement, avant le processus d'analyse, implique la définition des types de mots clés souhaités, pour l'indexation des séquences. En effet, les étiquettes associées aux objets au sein des images correspondent aux mots-clés qui seront prédits par le modèle lors de l’analyse. Dans le cadre du corpus étudié, il est pertinent d’extraire les mots-clés déjà indexés sur les \gls{sequence} au sein du logiciel. Ces mots-clés varient selon le contenu de l’image et peuvent, par exemple, être les suivants : affiche, foule, policier, gendarme, drapeau, etc. Ces mots-clés, tout comme le jeu de données COCO, sont assez génériques. De fait, le pré traitement est avant tout utile pour l’entraînement du modèle à pouvoir détecter des objets au sein d’images dont la qualité est compromise.   